\section{BACQ Benchmark}\label{bacq-benchmark}

Implements the framework of the BACQ Benchmark suite in a plateform
independent way. At the moment, it provides supports for the following
QPUs:

\begin{itemize}
\tightlist
\item
  Generic CPU (for reference)
\item
  Dwave
\item
  (TBA) Pasqal (Ruby machine from Genci)
\end{itemize}

And the following scores:

\begin{itemize}
\tightlist
\item
  G-Score (Evaluation of Quantum Optimization heuristics)
\item
  Lin2Score (Evaluation of capability for linear system solving in
  binary)
\end{itemize}

\subsection{📋 Prerequisites}\label{prerequisites}

\begin{itemize}
\tightlist
\item
  Python \textgreater= 3.10
\item
  CMake \textgreater= 3.18
\item
  Compilateur C++ (GCC/Clang)
\item
  Doxygen
\item
  Compilateur Latex (pdflatex)
\item
  Pandoc
\end{itemize}

\subsection{🚀 Installation}\label{installation}

\begin{itemize}
\tightlist
\item
  Build \& install
\end{itemize}

\begin{lstlisting}[language=bash]
git clone [https://github.com/utilisateur/mon-projet.git](https://github.com/utilisateur/mon-projet.git)
cd BacqBench
mkdir build && cd build
cmake ..
make
\end{lstlisting}

\begin{itemize}
\tightlist
\item
  Installation as python package (necessary for use with QPUs)

  \begin{itemize}
  \tightlist
  \item
    inside the main directory ``cd ..'' if required
  \item
    activate your current QPU environement (e.g.~DWave Ocean, Pasqal
    Pulser)
  \item
    install benchmark package using ``pip install -e .''
  \item
    N.B. you can do it after runing the tests (see below) Installing in
    the support environement of your QPU is necessary to run the QPU
    backend of the benchmark suite.
  \end{itemize}
\end{itemize}

\subsection{Basic unitary tests}\label{basic-unitary-tests}

\begin{lstlisting}[language=bash]
make check 
\end{lstlisting}

This provides several unitary tests. Not all are implemented yet.

\subsection{Usage}\label{usage}

The main python script is BACQBenchmark at the root of the project The
minimum parameter is the target QPU (or CPU comparison) E.g.

\begin{lstlisting}[language=bash]
python3 BACQBenchmak dwave 
\end{lstlisting}

Will launch the G-Score evaluation on the default architecture of DWave
(Advantage/Pegasus ATM). It is worth noting that this requires the
BACQBenchmak to be installed as package in the QPU environement. If it
is not the case, the main program would display ``Backend XXX not
available''

E.g.

\begin{lstlisting}[language=bash]
python3 BACQBenchmak cpu 
\end{lstlisting}

To run the G-Score on Simulated Annealing (CPU backend)

Use

\begin{lstlisting}[language=bash]
python3 BACQBenchmak -h
\end{lstlisting}

or

\begin{lstlisting}[language=bash]
python3 BACQBenchmak --help 
\end{lstlisting}

To be provided with usage. Each backend should provide its own help, but
in this case it is necessary to provide the backend on the command line

Being Alpha release, all options may not be implemented yet

\subsection{To do}\label{to-do}

Here is a preliminary of pending functionalities:

\begin{itemize}
\tightlist
\item
  Saving log in json
\item
  Evaluation of benchmark from json results file
\end{itemize}
